\documentclass[11pt,a4paper]{article}
\usepackage[utf8]{inputenc}
\usepackage[T1]{fontenc}
\usepackage[portuguese]{babel}
\usepackage{xcolor}
\usepackage{hyperref}
\usepackage{geometry}
\usepackage{tcolorbox}

\geometry{margin=2.5cm}

\definecolor{primary}{RGB}{39, 174, 96}
\definecolor{secondary}{RGB}{44, 62, 80}
\definecolor{codebg}{RGB}{240, 240, 240}

\hypersetup{colorlinks=true, linkcolor=primary, urlcolor=primary}

\title{\color{secondary}\Huge\textbf{Solução: Exercício 4.10 - O projeto, o grand finale}}
\author{\color{primary}DevOps com Kubernetes -- 2026}
\date{}

\begin{document}
\maketitle

\section*{\color{primary}Guia de Solução}
\begin{tcolorbox}[colback=green!5!white,colframe=green!60!black,title=Resolução Passo a Passo]
### Passo a Passo

1. \textbf{Separação de Repositórios}: Crie um repositório chamado \texttt{meu-projeto-codigo} (com arquivos Node/Go/Python) e um repositório \texttt{meu-projeto-manifestos} (apenas arquivos YAML/Kustomize).
2. \textbf{Pipeline de CI (Código)}: Configure o GitHub Actions no repositório de código para compilar as imagens e enviá-las ao Docker Hub, anexando o \texttt{\$GITHUB_SHA} como tag da imagem.
3. \textbf{Integração CI -> Manifestos}: Adicione um passo final ao GitHub Actions que faz um \texttt{git clone} do repositório \texttt{meu-projeto-manifestos} com um Token de Acesso Pessoal (PAT). Nesse diretório, rode \texttt{kustomize edit set image imagem-backend=meu-usuario/backend:\$GITHUB_SHA}, e faça o \textit{commit} e \textit{push} da mudança no repositório de manifestos.
4. \textbf{Pipeline de CD (ArgoCD)}: Configure o ArgoCD para escutar o repositório \texttt{meu-projeto-manifestos}. A alteração do \textit{commit} de CI acionará o \textit{Sync} automático do ArgoCD.
\end{tcolorbox}

\end{document}
